 VPF10, Amsterdam, NL.
\item[\textbullet] (Nov.~24) \emph{Dissipative anomaly in sliding drops}, APS-DFD, Salt Lake City, USA.
\item[\textbullet] (Sep.~24) \emph{Drop impact forces}, 12th Liquid Matter Conf., Mainz, DE.
\item[\textbullet] (Sep.~24) \emph{A unifying approach for drop impact dynamics on rigid surfaces}, 1st EFDC, Aachen, DE.
\item[\textbullet] (Apr.~24) \emph{Bursting bubbles in a viscoelastic medium}, European Rheology Conf., Leeds, UK.
\item[\textbullet] (Nov.~23) \emph{A unifying approach for droplet impact forces}, APS-DFD, Washington, DC, USA. \href{https://youtu.be/3S_hwkMCx6U?si=01tvC4RHCoga62bC}{\color{red}\faYoutube}
\item[\textbullet] (Nov.~22) \emph{Impact forces of water drops}, APS-DFD, Indianapolis, USA. \href{https://youtu.be/U0g67zV70hw}{\color{red}\faYoutube}
\item[\textbullet] (Jul.~23) \emph{Viscous free-surface flows}, Basilisk/Gerris Meeting, Paris, FR.
\item[\textbullet] (Sep.~22) \emph{When does an impacting drop stop bouncing?}, EFMC14, Athens, GR.
\item[\textbullet] (Jan.~22) \emph{How much force is required to play ping-pong with water droplets?}, Physics@Veldhoven, NL. \href{https://youtu.be/wmEYeQRPi0Q?si=u8Ex64dYLHQ8TxCE}{\color{red}\faYoutube}
\item[\textbullet] (Nov.~21) \emph{Viscous dissipation dictates Taylor-Culick type retractions}, APS-DFD, Phoenix, USA. \href{https://youtu.be/d0tMlpbMIqY?si=V-7Lh5g-zj1j7al6}{\color{red}\faYoutube}
\item[\textbullet] (Dec.~20) \emph{Bursting bubble in a viscoplastic medium}, International Congress on Rheology (Virtual). \href{https://youtu.be/52D6nfaHbYg?si=lE3oQGiMZe3UCT6j}{\color{red}\faYoutube}
\item[\textbullet] (Nov.~20) \emph{When does a viscous drop stop bouncing?}, APS-DFD (Virtual). \href{https://youtu.be/9JCV9JFwr2A?si=Vs1ljwdmZrzHUOXX}{\color{red}\faYoutube}
\item[\textbullet] (Feb.~20) \emph{Jumping \& Bouncing Drops \& Bubbles}, Max Planck meeting, Mainz, DE.
\item[\textbullet] (Nov.~19) \emph{Droplet Encapsulation}, APS-DFD, Seattle, USA.
\item[\textbullet] (Sep.~19) \emph{Bursting Bubbles: from Champagne to Mudpots}, VPF8, Cambridge, UK.
\item[\textbullet] (Aug.~19) \emph{Impinging drop lifts a sessile drop}, 9th 4U Summer School, DK.
\item[\textbullet] (May~16) \emph{On gas-liquid entrainment by impinging jet}, ICMF9, Florence, IT.
\end{itemize}

%===============================
%       SUMMARY / METRICS
%===============================
\section{Summary of Key Numbers (as of \today)}
\begin{itemize}[leftmargin=1.5em]
\item \faIdCard\hspace{0.3em}\textbf{Researcher ID:} \href{https://www.webofscience.com/wos/author/record/K-1856-2019}{K-1856-2019}
\item \faOrcid\hspace{0.3em}\textbf{Orcid:} \href{https://orcid.org/0000-0002-4293-6099}{0000-0002-4293-6099}
\item \faChartLine\hspace{0.3em}\textbf{Hirsch-index:} H = 11 (\href{https://scholar.google.com/citations?hl=en&user=tHb_qZoAAAAJ}{Google Scholar}), 9 (\href{https://www.webofscience.com/wos/author/record/K-1856-2019}{Web of Science})
\item \faFile\hspace{0.3em}\textbf{i10-index:} 12 (\href{https://scholar.google.com/citations?hl=en&user=tHb_qZoAAAAJ}{Google Scholar})
\item \faUsers\hspace{0.3em}\textbf{Research Interest Score:} \(1100+\) (top 2\% among \href{https://www.researchgate.net/profile/Vatsal-Sanjay-2}{ResearchGate} members who first published in 2015.)
\end{itemize}

% %------------------------------------------------------------------------------
% % POSSIBLE SUGGESTION:
% %  For a truly short CV (2-3 pages), you can move "Older contributed talks" or
% %  "Full publications list" to your personal website (and keep just a selection).
% %------------------------------------------------------------------------------

\end{document}
